\chapter{Pain relief after surgery}

\section*{Definition and Overview}
Postoperative pain management involves the prevention, evaluation, and treatment of pain in patients recovering from surgical procedures. Effective postoperative pain control is a critical component of perioperative care, essential not only for patient comfort and psychological well-being but also for mitigating the physiological stress response to surgery. Inadequate pain management can lead to significant systemic complications, including pulmonary dysfunction, cardiovascular stress, delayed wound healing, and prolonged hospital stays. Modern clinical practice advocates for a multimodal analgesia approach, which combines different classes of pain-relief medications and regional anaesthetic techniques to target multiple nociceptive pathways, thereby optimizing pain control while minimizing the side effects associated with high-dose opioid therapy.

\section*{Epidemiology}
Postoperative pain is a near-universal clinical challenge, affecting millions of patients undergoing surgical procedures annually. Despite improvements in anaesthetic techniques, studies indicate that approximately 80\% of surgical patients experience acute postoperative pain, with up to 40\% reporting their pain as moderate to severe. Major surgical procedures—such as thoracotomies, sternotomies, spinal surgeries, and major orthopaedic joint reconstructions (e.g., total knee or hip arthroplasty)—are consistently associated with the highest pain intensities. Furthermore, poorly controlled acute postoperative pain is the primary modifiable risk factor for the development of chronic postsurgical pain (CPSP), a debilitating condition that occurs in 10\% to 50\% of patients depending on the surgical site and technique.

\section*{Aetiology and Pathophysiology}
The pathophysiology of postoperative pain is complex, involving inflammatory, nociceptive, and neuropathic mechanisms initiated by tissue trauma.

\begin{itemize}
    \item \textbf{Inflammatory nociception:} Surgical incision, dissection, and tissue retraction cause direct cellular damage, triggering the release of inflammatory mediators (prostaglandins, bradykinin, histamine, cytokines, and substance P) into the surrounding tissue. These substances sensitise peripheral nociceptive terminals (primary hyperalgesia), lowering their threshold for firing.
    \item \textbf{Neuropathic injury:} Direct mechanical trauma, transection, stretching, or compression of peripheral nerves during the surgical procedure leads to ectopic neuronal firing and neuropathic pain symptoms.
    \item \textbf{Central sensitisation:} Continuous and intense nociceptive input from the wound site alters sensory processing in the dorsal horn of the spinal cord. This results in the "wind-up" phenomenon, characterized by hyperalgesia (increased response to painful stimuli) and allodynia (pain response to normally non-painful stimuli like light touch).
    \item \textbf{Visceral nociception:} In abdominal or pelvic surgeries, traction on the mesentery, distension of hollow organs, or peritoneal inflammation activates visceral nociceptors, producing a diffuse, poorly localised, aching or cramping pain.
    \item \textbf{Neuroendocrine stress response:} Acute pain stimulates the hypothalamic-pituitary-adrenal axis and the sympathetic nervous system, leading to elevated levels of cortisol, catecholamines, and aldosterone, which can cause systemic vasoconstriction and myocardial stress.
\end{itemize}

\section*{Clinical Features}
The clinical features of postoperative pain vary depending on the site of surgery, the type of anaesthesia used, and patient-specific variables.

\begin{itemize}
    \item \textbf{Localised incisional pain:} A sharp, stabbing, or burning pain localized to the surgical wound, characteristically exacerbated by physical movement, deep breathing, coughing, or wound palpation.
    \item \textbf{Deep visceral aching:} A dull, generalised, or cramping pain in the abdomen or pelvis, often associated with abdominal distension, nausea, and delayed passage of flatus or stool.
    \item \textbf{Referred pain:} Pain perceived in a site distant from the surgical wound, such as shoulder tip pain following laparoscopic surgery due to diaphragmatic irritation from carbon dioxide gas insufflation.
    \item \textbf{Skeletal muscle spasm:} Involuntary, painful contraction of muscles surrounding the surgical site, common in spinal and orthopaedic procedures.
    \item \textbf{Sympathetic hyperactivity:} Physiological signs of pain, including tachycardia, hypertension, tachypnoea, diaphoresis, and pupillary dilation.
    \item \textbf{Functional restriction:} Guarding behaviour, shallow respirations, reluctance to mobilise or participate in chest physiotherapy, and severe sleep fragmentation.
\end{itemize}

\section*{Differential Diagnosis}
It is critical to distinguish typical, expected postoperative pain from pain indicating acute, life-threatening surgical complications or systemic events.

\begin{itemize}
    \item \textbf{Postoperative haemorrhage:} Presents as sudden, severe, localized pain, escalating swelling, abdominal distension, signs of hypovolaemic shock (hypotension, tachycardia, pallor), and a rapid decline in haemoglobin concentration.
    \item \textbf{Surgical site infection (SSI):} Typically manifests 3 to 7 days postoperatively with worsening, localized wound pain, erythema, localized warmth, induration, purulent discharge, and systemic pyrexia.
    \item \textbf{Anastomotic leak or visceral perforation:} Following gastrointestinal surgery, this presents with severe, generalised abdominal pain, abdominal rigidity, guarding, rebound tenderness (peritonitis), high fever, and septic shock.
    \item \textbf{Deep vein thrombosis (DVT) or pulmonary embolism (PE):} Presents as unilateral calf or thigh pain, swelling, warmth, and erythema (DVT), or sudden-onset pleuritic chest pain, dyspnoea, tachypnoea, and hypoxaemia (PE).
    \item \textbf{Compartment syndrome:} A surgical emergency occurring in limbs, characterized by severe, progressive pain out of proportion to clinical findings, tense swelling, pain on passive stretch, paraesthesia, and late-stage pulselessness.
\end{itemize}

\section*{When to Seek Medical Care}
Patients recovering at home or in hospital wards should alert their medical team immediately if they experience pain that is not controlled by their prescribed medications or if they note new symptoms.

\textbf{Call 000 (or local emergency services) for:}
\begin{itemize}
    \item Sudden, severe shortness of breath, rapid breathing, or pleuritic chest pain.
    \item Rapid soaking of surgical dressings with bright red blood, or signs of internal bleeding.
    \item Unresponsiveness, severe somnolence, pinpoint pupils, or slowed breathing (indicating opioid toxicity).
    \item Severe, generalised abdominal pain with abdominal rigidity, high fever, and vomiting.
    \item Loss of motor function, numbness, or loss of bowel or bladder control in patients who have recently had an epidural or spinal catheter.
\end{itemize}

\section*{Diagnosis and Investigations}
The diagnosis of postoperative pain relies on systematic clinical assessment combined with investigations to rule out underlying surgical pathology.

\begin{itemize}
    \item \textbf{Subjective pain scores:} Routine evaluation of pain intensity using validated tools such as the Numerical Rating Scale (NRS) or Visual Analogue Scale (VAS) at rest and during functional activities (e.g., coughing, mobilising).
    \item \textbf{Clinical wound assessment:} Inspection of the surgical dressing and wound site for signs of haematoma, dehiscence, erythema, warmth, or purulent fluid collection.
    \item \textbf{Haematological and biochemical profiles:} Full blood count (to monitor for anaemia or leukocytosis), inflammatory markers (C-reactive protein), and serum electrolytes.
    \item \textbf{Radiological investigations:} Chest X-rays, Doppler ultrasound for suspected DVT, or CT scans of the abdomen/pelvis with contrast to evaluate for anastomotic leaks or fluid collections.
    \item \textbf{Cardiorespiratory monitoring:} Pulse oximetry, capnography, and respiratory rate monitoring to screen for hypoxaemia and respiratory depression secondary to systemic analgesics.
\end{itemize}

\section*{Management}
Postoperative pain management is highly individualised, utilizing a multimodal strategy to target multiple pain pathways and minimize adverse drug events.

\begin{itemize}
    \item \textbf{Simple non-opioid analgesics:} Scheduled administration of paracetamol and non-steroidal anti-inflammatory drugs (NSAIDs, e.g., ibuprofen, ketorolac) or selective COX-2 inhibitors (e.g., celecoxib), which form the foundation of postoperative pain therapy.
    \item \textbf{Patient-Controlled Analgesia (PCA):} Intravenous infusion of strong opioids (e.g., morphine, fentanyl, hydromorphone) self-administered by the patient via a programmed pump with a lock-out interval to prevent overdose.
    \item \textbf{Regional and neuraxial anaesthesia:} Continuous epidural infusions of local anaesthetics and opioids, or single-injection/catheter-guided peripheral nerve blocks (e.g., femoral nerve block for knee replacement; transversus abdominis plane [TAP] block for abdominal incisions).
    \item \textbf{Adjuvant medications:} Low-dose ketamine infusions (NMDA receptor antagonist), gabapentinoids (gabapentin or pregabalin) for neuropathic pain components, and intravenous lidocaine infusions to reduce central sensitisation and opioid requirements.
    \item \textbf{Non-pharmacological and physical therapies:} Splinting the surgical incision with a pillow during coughing or deep breathing, thermotherapy (ice or heat packs), relaxation techniques, and early, structured physiotherapy.
\end{itemize}

\section*{Complications}
The interventions and medications used to treat postoperative pain can cause significant complications that must be actively monitored and managed.

\begin{itemize}
    \item \textbf{Opioid-induced respiratory depression:} The most critical complication, characterized by a decrease in respiratory rate and tidal volume, leading to hypercapnia, hypoxia, and potential respiratory arrest.
    \item \textbf{Opioid-induced bowel dysfunction:} Severe constipation, abdominal distension, and postoperative ileus due to the activation of mu-opioid receptors in the enteric nervous system, which delays gut motility.
    \item \textbf{Postoperative nausea and vomiting (PONV):} A common side effect of systemic opioids and anaesthetic agents, causing patient distress, dehydration, and potential wound disruption.
    \item \textbf{Neuraxial complications:} Spinal epidural haematoma or abscess, post-dural puncture headache, and urinary retention secondary to block-induced bladder detrusor hypocontractility.
    \item \textbf{Adverse drug toxicity:} NSAID-induced gastrointestinal ulceration, acute kidney injury (particularly in hypovolaemic patients), or paracetamol-induced hepatotoxicity in cases of accidental overdose.
\end{itemize}

\section*{Prognosis and Outcomes}
With a well-structured multimodal analgesia protocol, the prognosis for acute postoperative pain relief is excellent. Pain typically peaks within the first 24 to 72 hours postoperatively and declines progressively as tissue healing occurs. Effective acute pain management is associated with faster recovery of pulmonary function, earlier ambulation, reduced cardiovascular events, and shorter hospital stays. Most patients successfully transition to oral analgesics within a few days and discontinue pain medications entirely within several weeks. Furthermore, optimized acute pain control significantly decreases the risk of transitioning from acute postoperative pain to chronic postsurgical pain syndromes.

\section*{Prevention and Self-Care}
Preoperative planning and post-discharge self-care are essential for optimizing pain management and recovery at home.

\begin{itemize}
    \item \textbf{Preoperative education:} Discussing pain expectations, the multimodal pain plan, and instructions on how to use PCA devices with the surgical and anaesthetic teams.
    \item \textbf{Strict adherence to medication schedules:} Taking non-opioid medications (paracetamol and NSAIDs) on a regular schedule rather than waiting for severe pain to return, to maintain therapeutic drug levels.
    \item \textbf{Incision splinting:} Supporting the surgical wound with a small pillow or hand when coughing, sneezing, deep breathing, or changing positions to minimise mechanical tension and pain.
    \item \textbf{Systematic opioid tapering:} Reducing the dose and frequency of opioid medications as the pain improves, transitioning to simple analgesics, and safely disposing of unused opioids.
    \item \textbf{Early and safe mobility:} Engaging in gentle physical activities (e.g., walking, ankle pumps) as prescribed to prevent thromboembolic events, promote bowel motility, and reduce joint stiffness.
\end{itemize}

\section*{Sources and Further Reading}
The following reputable organisations provide further information on postoperative pain management:

\begin{itemize}
    \item Australian and New Zealand College of Anaesthetists (ANZCA). \textit{Acute Pain Management: Scientific Evidence.} https://www.anzca.edu.au/
    \item Healthdirect Australia. \textit{Managing pain after surgery.} https://www.healthdirect.gov.au/
    \item American Society of Anesthesiologists (ASA). \textit{Managing Pain After Surgery.} https://www.asahq.org/
    \item National Institute for Health and Care Excellence (NICE). \textit{Acute pain management after surgery.} https://www.nice.org.uk/
    \item World Health Organization (WHO). \textit{WHO guidelines for the pharmacological and radiotherapeutic management of cancer pain in adults and adolescents.} https://www.who.int/
\end{itemize}
